% Reusable finite bound; optional research-note material, not needed for K4 exclusion.
\subsection{反対符号の二つの例外 reverse が大きい場合の有限化}
\label{fc17:tail-note}

低 row と column が
\[
 r=(P,-P+h-1,-h),\qquad x=(1,-1,-v),\qquad
 P\ge h\ge1,\quad v\ge1
\]
である整数三高分母模型を考える．分母は
\[
 q=(P+1,P-h+2,hv+1)
\]
となる．高高 entry $t=p_{13}$ を用いて他の entries を行和・pair式から消去すると，
\[
 P\{vt^2+2(hv+1)t+(h-1)(hv+1)\}
 +vt^2+\{v(h+1)+1-h\}t+hv(h-1)=0
\]
を得る．$t\ne0$ であり，$t>0$ なら左辺は正なので，$t=-z$，$z\ge1$ でなければならない．
\[
 \begin{split}
 A(z)&=vz^2-2(hv+1)z+(h-1)(hv+1),\\
 B(z)&=vz^2-\{v(h+1)+1-h\}z+hv(h-1)
 \end{split}
\]
と置けば $PA(z)+B(z)=0$ であり，
\[
 B(z)-A(z)=\{(h-1)v+h+1\}z-(h-1)>0.
\]
従って $A(z)<0$ で，
\[
 P+1=\frac{\{(h-1)v+h+1\}z-(h-1)}{-A(z)}.
\]
$A(z)=B(z)=0$ という共鳴は生じない．

$z\ge2h+2$ なら
\[
 A(2h+2)=(v-1)(h^2+3h+4)+(h^2-1)\ge0
\]
であり，その先で $A$ は増加する．また
\[
 B(z)=vz(z-h-1)+(h-1)z+hv(h-1)>0
\]
なので矛盾する．従って
\[
 1\le z\le2h+1,\qquad
 P\le\{(h-1)v+h+1\}(2h+1)-h.
\]
よって $h,v$ を固定すれば全 $P$ を完全に復元できる．
全分母が奇数なら $P$ は偶数，$h$ は奇数である．
通常 reverse の上界 $h\le4$ が使える場合には $h=1,3$ のみとなる．
この有限化は他の graph で反対符号の大きい例外 entries が残る場合に転用できる．
